\section{基础算法}




%\subsection{高精度计算}


\subsection{二分查找算法}

\begin{lstlisting}[language=c++,caption={查找最后一个小于等于q的数的下标}]
//查找最后一个小于等于q的数的下标
int find(int q) {
	int l = 0, r = n + 1; //开区间
	while (l + 1 < r) {
		// l + 1 < r时结束
		int mid = l + r >> 1;
		if (a[mid] <= q) {
			l = mid;
		} else {
			r = mid;
		}
	}
	return l;
}
\end{lstlisting}

\begin{lstlisting}[language=c++,caption={查找第一个大于等于q的数的下标}]
//查找第一个大于等于q的数的下标
int find(int q) {
	int l = 0, r = n + 1; //开区间
	while (l + 1 < r) {
		// l + 1 < r时结束
		int mid = l + r >> 1;
		if (a[mid] >= q) {
			r = mid;
		} else {
			l = mid;
		}
	}
	return r;
}

\end{lstlisting}

\begin{lstlisting}[language=c++,caption={洛谷P2249 查找}]
//洛谷P2249 【深基13.例1】查找
//询问数字,输出这个数字在序列中第一次出现的编号，如果没有找到的话输出 −1 
#include<bits/stdc++.h>
using namespace std;
int main() {
	ios::sync_with_stdio(false);
	cin.tie(nullptr);
	int n, m;
	cin >> n >> m;
	vector<int> v(n + 2, 0);
	for (int i = 1; i <= n; i++) {
		cin >> v[i];
	}
	while (m--) {
		int q;
		cin >> q;
		auto find = [&](const int k) {
			int l = 0, r = n + 1;
			while (l + 1 < r) {
				int mid = (l + r) >> 1;
				if (v[mid] >= k) {
					r = mid;
				} else {
					l = mid;
				}
			}
			return r;
		};
		if (v[find(q)] == q) {
			cout << find(q) << " ";
		} else {
			cout << "-1" << " ";
		}
	}
	return 0;
}

\end{lstlisting}


\begin{lstlisting}[language=c++,caption={洛谷P1024 一元三次方程求解}]
//洛谷P1024 [NOIP 2001 提高组]一元三次方程求解
#include<bits/stdc++.h>
using namespace std;
double a, b, c, d;

int main() {
	cin >> a >> b >> c >> d;
	auto func = [&](double x)->double {
		return a * x * x * x + b * x * x + c * x + d;
	};
	auto find = [&](double l, double r)->double {
		while (r - l > 0.0001) {
			double mid = (l + r) / 2;
			if (func(mid) * func(r) < 0) {
				l = mid;
			}else {
				r = mid;
			}
		}
		return l;
	};
	for (int i = -100; i <= 100; i++) {
		if (func(i) == 0) {
			cout << fixed << setprecision(2) << (double)i << " ";
		}
		if (func(i) * func(i + 1) < 0) {
			cout << fixed << setprecision(2) << (double)find(i, i + 1) << " ";
		}
	}
}
\end{lstlisting}

\subsection{二分答案算法}


\subsection{分数规划}


\subsection{前缀和 \quad 二维前缀和}


\subsection{树上前缀和}

\subsection{差分 \quad 二维差分}

\subsection{树上差分}

\subsection{ST表 \quad RMQ问题}

\subsection{快速排序 \quad 第k小的数}

\subsection{归并排序 \quad 逆序对}

\subsection{堆 \quad 堆排序}

\subsection{双指针}

\subsection{区间合并}

\subsection{贪心算法}